\exercisesection{}

In the exercises of this section, unless otherwise stated, G denotes a locally compact topological group, equipped with a left Haar measure which is denoted by \(\mu\) .

\begin{enumerate}
\def\labelenumi{\arabic{enumi})}
\item
  Let \(\pi\) be a unitary representation of G on a complex Hilbert space E. The representation \(\pi\) is irreducible if and only if the image of \(\mathcal{H}(\mathrm{G})\) in \(\mathcal{L}(\mathrm{E})\) is dense for the topology of pointwise convergence. (Use Exercise 24 of IV, p. 326.)
\item
  Let H be a closed subgroup of G. Let \(\nu\) be a quasi-invariant measure on G/H (INT, VII, p. 56, § 2, n° 5, th. 2), and let \(\varrho\) be the function on G with values in \(\mathbf{R}_{+}^{*}\) associated with it. The maps \(\varrho(g)\) defined by
\end{enumerate}

\[
\varrho (g) f (x) = \left(\frac {\varrho (g ^ {- 1} x)}{\varrho (x)}\right) ^ {1 / 2} f (g ^ {- 1} x)
\]

define a continuous linear representation of G on \(\mathcal{L}^{2}(\mathrm{G / H},\nu)\) , and, by passing to the quotient, define a unitary representation of G on \(\mathrm{L}^2 (\mathrm{G / H},\nu)\) , called the quasi-regular representation of \(\mathrm{G / H}\) .

\begin{enumerate}
\def\labelenumi{\arabic{enumi})}
\setcounter{enumi}{2}
\tightlist
\item
  Let G = R and let H be the group R equipped with the discrete topology and with the counting measure, which is a Haar measure.
\end{enumerate}

\begin{enumerate}
\def\labelenumi{\alph{enumi})}
\item
  For every \(t \in G\) , and every \(f \in L^{2}(H)\) , the map \(\varrho(t)f: x \mapsto f(x + t)\) belongs to \(L^{2}(H)\) ; the map \(\varrho(t)\) is a unitary map on \(L^{2}(H)\) . b) For all \(f,g \in L^{2}(H)\) , the map from G to C defined by \(t \mapsto \langle f | \varrho(t)g \rangle\) is measurable.
\item
  Let \(f\) and \(g\) be \(\mathrm{L}^2(\mathrm{H})\) . If there exists \(t \in \mathbf{R}\) such that \(\langle f | \varrho(t)g \rangle\) is nonzero, then the map \(t \mapsto \langle f | \varrho(t)g \rangle\) from \(\mathbf{G}\) into \(\mathbf{C}\) is not continuous. (In particular, the hypothesis that E is countably generated cannot be omitted from cor. 2 of INT, VIII, p. 171, §4, n° 7.)
\end{enumerate}

\(d)\) There exists \(f \in \mathrm{L}^2(\mathrm{H})\) such that the subspace of \(\mathrm{L}^2(\mathrm{H})\) generated by the elements \(\varrho(t)f\) , where \(t \in \mathbf{R}\) , is dense in \(\mathrm{L}^2(\mathrm{H})\) .

\begin{enumerate}
\def\labelenumi{\arabic{enumi})}
\setcounter{enumi}{3}
\item
  Let \(p \geqslant 1\) . The kernel of the biregular representation of G on \(\mathrm{L}^{p}(\mathrm{G}, \mu)\) is the image of the center of G under the homomorphism \(g \mapsto (g, g)\) .
\item
  Let G be a locally compact unimodular group.
\end{enumerate}

\begin{enumerate}
\def\labelenumi{\alph{enumi})}
\item
  If there exists a finite-dimensional square-integrable representation of G, then G is compact.
\item
  If there exists a square-integrable representation of G, then the center of G is compact.
\end{enumerate}

\begin{enumerate}
\def\labelenumi{\arabic{enumi})}
\setcounter{enumi}{5}
\tightlist
\item
  We assume that the group G is unimodular. We denote by \(\widehat{G}_{d} \subset \widehat{G}\) the set of classes of irreducible unitary square-integrable representations of G. For \(\pi \in \widehat{G}_{d}\) , we denote by \(p_{\pi}\) the orthogonal projector on \(L^{2}(G)\) whose image is the \(\pi\)-isotypic component of the right regular representation of G.
\end{enumerate}

\(a)\) Let \(f \in \mathcal{L}^2(\mathrm{G})\) and let \(\pi \in \widehat{\mathrm{G}}_d\) . We denote by E the representation space of \(\pi\) . The map \(g \mapsto f(g)\pi(g)\) belongs to \(\mathrm{L}^1(\mathrm{G};\mathcal{L}(\mathrm{E}))\) ; the endomorphism

\[
v = \int_ {\mathrm{G}} f (g) \pi (g) d \mu (g)
\]

is a Hilbert--Schmidt endomorphism, denoted \(v = \pi(f)\) .

\begin{enumerate}
\def\labelenumi{\alph{enumi})}
\setcounter{enumi}{1}
\item
  Let \(f \in \mathcal{L}^2(\mathrm{G})\) . We have \(\| \pi(\overline{f}) \|_2 = \| p_\pi(f) \|\) .
\item
  Let \(L_{d}^{2}(G)\) be the closed subspace of \(L^{2}(G)\) generated by the \(\pi\)-isotypic components of \(L^{2}(G)\) for \(\pi \in \widehat{G}_{d}\) . It is the Hilbert sum of the images of the orthogonal projectors \(p_{\pi}\) for \(\pi \in \widehat{G}_{d}\) . For every \(f \in L_{d}^{2}(G)\) , we have
\end{enumerate}

\[
\| f \| ^ {2} = \sum_ {\pi \in \widehat {\mathrm{G}} _ {d}} c (\pi) \| \pi (f) \| _ {2} ^ {2}
\]

where \(c(\pi)\) is the formal degree of \(\pi\) relative to \(\{e\}\) .

\begin{enumerate}
\def\labelenumi{\arabic{enumi})}
\setcounter{enumi}{6}
\tightlist
\item
  Let \(n \in N\) . Determine the unitary automorphic representations of \(R^{n}\) (that is, the representations \(\Gamma\) -automorphic for every discrete subgroup \(\Gamma\) of \(R^{n}\) ).
\end{enumerate}

¶ 8) We assume that G is a locally compact unimodular group. Let Γ be a discrete subgroup of G such that the quotient space X = Γ\textbackslash G is compact. Let μ be a Haar measure on G and β the counting measure on Γ; we denote by \(\widetilde{\mu} = \mu/\beta\) .

Let \(\varrho\) be the unitary representation of G on \(\mathrm{L}^{2}(\mathrm{X},\widetilde{\mu})\) defined in Example 3 of V, p. 404.

For \(\gamma\in\Gamma\) , we denote by \(\Gamma_{\gamma}\) (resp. \(\mathbf{G}_{\gamma}\) ) the centralizer of \(\gamma\) in \(\Gamma\) (resp. \(\mathbf{G}\) ). \(a)\) The group \(\mathbf{G}_{\gamma}\) is unimodular.

\begin{enumerate}
\def\labelenumi{\alph{enumi})}
\setcounter{enumi}{1}
\item
  Let \(\varphi\in\mathcal{K}(\mathrm{G})\) . For every \(\Gamma\)-automorphic unitary representation \(\pi\) , the endomorphism \(\pi(\varphi)\) has finite trace.
\item
  One has
\end{enumerate}

\[
\operatorname{Tr} (\varrho (\varphi)) = \sum_ {\pi \in \widehat {\mathrm{G}}} m _ {\pi} (\varrho) \operatorname{Tr} (\pi (\varphi))
\]

where \(m_{\pi}(\varrho)\) is the multiplicity of \(\pi\) in \(\varrho\) (“spectral side of the trace formula”).

\begin{enumerate}
\def\labelenumi{\alph{enumi})}
\setcounter{enumi}{3}
\tightlist
\item
  Let \(\Gamma^{\sharp}\) be the set of conjugacy classes of \(\Gamma\) . For \(\gamma \in \Gamma^{\sharp}\) , we denote by \(\beta_{\gamma}\) the counting measure on \(\Gamma_{\gamma}\) . We have
\end{enumerate}

\[
\mathrm{Tr} (\varrho (\varphi)) = \sum_ {\gamma \in \Gamma^ {\sharp}} \int_ {\Gamma_ {\gamma} \backslash G} \varphi (x ^ {- 1} \gamma x) d (\mu / \beta_ {\gamma}) (x).
\]

\begin{enumerate}
\def\labelenumi{\alph{enumi})}
\setcounter{enumi}{4}
\tightlist
\item
  We fix a Haar measure \(\mu_{\gamma}\) on \(G_{\gamma}\) for all \(\gamma \in \Gamma\) . We have
\end{enumerate}

\[
\operatorname{Tr} (\varrho (\varphi)) = \sum_ {\gamma \in \Gamma^ {\sharp}} \int_ {G _ {\gamma} \backslash G} \int_ {\Gamma_ {\gamma} \backslash G _ {\gamma}} \varphi (x ^ {- 1} y ^ {- 1} \gamma y x) d (\mu_ {\gamma} / \beta_ {\gamma}) (y) d (\mu / \mu_ {\gamma}) (x).
\]

\(f)\) For every \(\gamma \in \Gamma\) , the quotient measure \(\mu_{\gamma} / \beta_{\gamma}\) of \(\Gamma_{\gamma} \backslash G_{\gamma}\) is finite. \(g)\) One has

\[
\mathrm{Tr} (\varrho (\varphi)) = \sum_ {\gamma \in \Gamma^ {\sharp}} (\mu_ {\gamma} / \beta_ {\gamma}) (\Gamma_ {\gamma} \backslash \mathrm{G} _ {\gamma}) \int_ {\mathrm{G} _ {\gamma} \backslash \mathrm{G}} \varphi (x ^ {- 1} \gamma x) d (\mu / \mu_ {\gamma}) (x)
\]

(“geometric side of the trace formula”; the formula expressing the equality of this expression with that obtained in question c) is called the “trace formula” for Γ-automorphic representations).

\begin{enumerate}
\def\labelenumi{\alph{enumi})}
\setcounter{enumi}{7}
\tightlist
\item
  Set G = R and \(\Gamma = Z\) . Verify that the trace formula is equivalent in this case to the Poisson formula (cf. II, p. 239).
\end{enumerate}

\begin{enumerate}
\def\labelenumi{\arabic{enumi})}
\setcounter{enumi}{8}
\item
  We assume that G is a perfect real semisimple Lie group without compact factor. There is no nontrivial finite-dimensional unitary representation of G.
\item
  \begin{enumerate}
  \def\labelenumii{\alph{enumii})}
  \tightlist
  \item
    Every matrix coefficient of the left regular representation of G on \(L^{2}(G)\) belongs to \(\mathcal{C}_{0}(G)\) .
  \end{enumerate}
\end{enumerate}

\begin{enumerate}
\def\labelenumi{\alph{enumi})}
\setcounter{enumi}{1}
\tightlist
\item
  If G is not compact, the left regular representation of G on L²(G) contains no nonzero finite-dimensional subrepresentation.
\end{enumerate}

\begin{enumerate}
\def\labelenumi{\arabic{enumi})}
\setcounter{enumi}{10}
\tightlist
\item
  Let E be a Hilbert space.
\end{enumerate}

\begin{enumerate}
\def\labelenumi{\alph{enumi})}
\item
  Let \(\varrho\) be a unitary representation of the group \(\mathbf{R}\) in E. Let D be a dense subspace of E stable under \(\varrho\) such that, for every \(x\in\mathrm{D}\) , the map \(\psi_{x}:t\mapsto\varrho(t)x\) is differentiable on R. Let u be the partial operator with domain D such that \(u(x) = i^{-1}\psi_x'(0)\) for all \(x \in D\) . Then u is essentially self-adjoint and \(\overline{u}\) is the infinitesimal generator of \(\varrho\) .
\item
  Let u be a self-adjoint partial operator on E and let D be a dense subspace of \(\text{dom}(u)\) . Suppose that for every \(t \in R\) , we have \(\exp(itu)(D) \subset D\) . Then the space D is dense in \(E_{u}\) .
\item
  Let u be a self-adjoint partial operator on E. Let F be a Hilbert space and v a self-adjoint partial operator on F. The partial operator \(u \otimes 1_{F} + 1_{E} \otimes E\) on \(E \otimes F\) defines an essentially self-adjoint partial operator on \(E \widehat{\otimes}_{2} F\) .
\end{enumerate}

\begin{enumerate}
\def\labelenumi{\arabic{enumi})}
\setcounter{enumi}{11}
\tightlist
\item
  Let G be a real Lie group and let \(\varrho\) be a unitary representation of G in a Hilbert space E. We say that an element \(x \in \mathrm{E}\) is a vector of class \(\mathrm{C}^{\infty}\) if the map from G to E defined by \(g \mapsto \varrho(g)x\) is of class \(\mathrm{C}^{\infty}\) . a) Let \(\pi\) be a unitary representation of G in F and \(u: \mathrm{E} \to \mathrm{F}\) a G-morphism. For every vector \(x \in \mathrm{E}\) of class \(\mathrm{C}^{\infty}\) , the vector \(u(x)\) is of class \(\mathrm{C}^{\infty}\) in F.
\end{enumerate}

We denote by \(\mathrm{C}^{\infty}(\varrho)\) the space of vectors of class \(C^{\infty}\) in E. By question a), every G-morphism from \(\varrho\) into \(\pi\) defines, by passing to subspaces, a linear map from \(\mathrm{C}^{\infty}(\varrho)\) into \(\mathrm{C}^{\infty}(\pi)\) .

\begin{enumerate}
\def\labelenumi{\alph{enumi})}
\setcounter{enumi}{1}
\tightlist
\item
  The space \(\mathrm{C}^{\infty}(\varrho)\) is dense in E.
\end{enumerate}

\begin{enumerate}
\def\labelenumi{\arabic{enumi})}
\setcounter{enumi}{12}
\tightlist
\item
  Let \(\pi\) be an irreducible unitary representation of G in a complex Hilbert space E. Let \(\varrho\) be a unitary representation of G in a complex Hilbert space F. Suppose there exists a closed densely defined partial operator \(u\) from E into F whose domain is a G-stable subspace of E and which satisfies \(u \circ \pi(g) = \varrho(g) \circ u\) for all \(g \in \mathbf{G}\) .
\end{enumerate}

The partial operator \(u\) belongs to \(\mathcal{L}(\mathrm{E};\mathrm{F})\) , and there exists \(\lambda > 0\) such that \(\lambda u\) is isometric.

\begin{enumerate}
\def\labelenumi{\arabic{enumi})}
\setcounter{enumi}{13}
\tightlist
\item
  Let Z be a subgroup of the center C of G such that C/Z is compact and let \(\nu\) be a Haar measure on G/Z. Let \(\pi\) be an irreducible unitary representation of G in a complex Hilbert space E. We assume that \(\pi\) is square-integrable modulo the center. Let \(u \in \mathcal{L}(\mathrm{E})\) be a nuclear endomorphism (def. 4 of IV, p. 169).
\end{enumerate}

\begin{enumerate}
\def\labelenumi{\alph{enumi})}
\tightlist
\item
  For every \(x\) and \(y\) in E, the function
\end{enumerate}

\[
g \mapsto \langle x | \pi (g) u \pi (g ^ {- 1}) y \rangle
\]

belongs to \(\mathcal{F}(\mathrm{G / Z})\) .

\begin{enumerate}
\def\labelenumi{\alph{enumi})}
\setcounter{enumi}{1}
\tightlist
\item
  For every \(x\) and \(y\) in E, we have
\end{enumerate}

\[
\int_ {\mathrm{G} / \mathrm{Z}} \langle x \mid \pi (g) u \pi (g ^ {- 1}) y \rangle d \nu (g) = \frac {1}{c _ {\mathrm{Z}} (\pi)} \langle x \mid y \rangle \operatorname{Tr} (u).
\]

\begin{enumerate}
\def\labelenumi{\arabic{enumi})}
\setcounter{enumi}{14}
\tightlist
\item
  Let Z be a subgroup of the center C of G such that C/Z is compact and let \(\nu\) be a Haar measure on G/Z. Let \(\pi\) be an irreducible unitary representation of G in a complex Hilbert space E. Denote by \(\chi\) the restriction to Z of the central character of \(\pi\) .
\end{enumerate}

We denote by \(f_{x,y}\) the matrix coefficient \(g \mapsto \langle x | \pi(g)y \rangle\) for \((x, y) \in \mathrm{E}^{2}\) . Let \(p \geqslant 1\) be a real number and let \(q\) be the conjugate exponent of p.

\begin{enumerate}
\def\labelenumi{\alph{enumi})}
\item
  Let \(y \in \mathrm{E}\) such that the matrix coefficients \(f_{x,y}\) belong to \(\mathcal{L}_{\chi}^{p}(\mathrm{G})\) . The map \(x \mapsto f_{x,y}\) is a continuous G-morphism from \(\pi\) into \(\boldsymbol{\delta}_{\mathrm{G},\chi}^{(p)}\) .
\item
  Let \(\varrho\) be an irreducible unitary representation of G in a complex Hilbert space F whose central character restricts to \(\chi\) on \(Z\) . One assumes that \(p > 1\) and that the matrix coefficients of \(\pi\) belong to \(\mathcal{L}_{\chi}^{p}(\mathrm{G})\) and those of \(\varrho\) belong to \(\mathcal{L}_{\chi}^{q}(\mathrm{G})\) . Let \(x_0 \in \mathrm{E}\) and \(y_0 \in \mathrm{F}\) . For every \(x \in \mathrm{E}\) there exists a continuous linear form \(\ell_x\) from F to C such that
\end{enumerate}

\[
\ell_ {x} (y) = \int_ {\mathrm{G/Z}} \overline {{\langle x _ {0} | \pi (g) x \rangle}} \langle y _ {0} | \varrho (g) y \rangle d \nu (g)
\]

for all \(y \in \mathbf{F}\) .

\begin{enumerate}
\def\labelenumi{\alph{enumi})}
\setcounter{enumi}{2}
\item
  There exists a continuous linear map u from E into F such that \(\ell_{x}(y)=\langle u(x)\mid y\rangle\) for all \((x,y)\in\mathrm{E}\times\mathrm{F}\) .
\item
  One has
\end{enumerate}

\[
\int_ {\mathrm{G} / \mathrm{Z}} \overline {{\langle x | \pi (g) x ^ {\prime} \rangle}} \langle y | \varrho (g) y ^ {\prime} \rangle d \nu (g) = 0
\]

for \((x,x^{\prime},y,y^{\prime})\in \mathrm{E}^{2}\times \mathrm{F}^{2}\)

\begin{enumerate}
\def\labelenumi{\alph{enumi})}
\setcounter{enumi}{4}
\tightlist
\item
  We assume that p = 1. The matrix coefficients of \(\pi\) are orthogonal to those of \(\varrho\) .
\end{enumerate}

\begin{enumerate}
\def\labelenumi{\arabic{enumi})}
\setcounter{enumi}{15}
\tightlist
\item
  Let \(\mathbf{G} = \mathbf{S}\mathbf{L}(2, \mathbf{R})\) and let H be the open subset of C formed by the complex numbers z such that \(\mathcal{I}(z) > 0\) . Let \(\mu\) be the Lebesgue measure on C. For every element
\end{enumerate}

\[
g = \left( \begin{array}{c c} a & b \\ c & d \end{array} \right)
\]

of G and every \(z \in \mathbf{H}\) we set \(j(g, z) = cz + d\) .

\begin{enumerate}
\def\labelenumi{\alph{enumi})}
\tightlist
\item
  The group G acts continuously on the left on H by
\end{enumerate}

\[
\left( \begin{array}{c c} a & b \\ c & d \end{array} \right) \cdot z = \frac {a z + b}{c z + d}.
\]

The measure \(\nu = y^{-2} \cdot \mu\) is G-invariant.

\begin{enumerate}
\def\labelenumi{\alph{enumi})}
\setcounter{enumi}{1}
\tightlist
\item
  Let n be an integer such that \(n \geqslant 2\) . We denote by \(E_{n}\) the vector space of holomorphic functions \(f: H \to C\) such that the function \(z \mapsto |f(z)|^{2} \mathcal{I}(z)^{n}\) belongs to \(\mathcal{L}^{2}(\mathbf{H}, \nu)\) . The space \(E_{n}\) endowed with the sesquilinear form
\end{enumerate}

\[
\langle f \mid g \rangle = \int_ {\mathbf {H}} \overline {{f (z)}} g (z) \mathcal {I} (z) ^ {n} d \nu (z)
\]

is a Hilbert space of countable type and of infinite dimension.

\begin{enumerate}
\def\labelenumi{\alph{enumi})}
\setcounter{enumi}{2}
\tightlist
\item
  For every \(g \in G\) and every \(f \in E_{n}\) we set
\end{enumerate}

\[
\pi_ {n} (g) f (z) = j (g ^ {- 1}, z) ^ {- n} f (g ^ {- 1} \cdot z).
\]

The map \(\pi_n\) defines a unitary representation of G in \(\mathrm{E}_n\) . \(d)\) The function \(f_n\) defined by \(f_n(z) = (z + i)^{-n}\) belongs to \(\mathrm{E}_n\) ; it satisfies \(\pi_n(k)f_n = \chi_{-n}(k)f_n\) for all \(k \in \mathbf{SO}(\mathbf{R}^2)\) , where \(\chi_n\) is the character such that

\[
\chi_ {n} \Big (\left( \begin{array}{c c} \cos (\theta) & \sin (\theta) \\ - \sin (\theta) & \cos (\theta) \end{array} \right) \Big) = e ^ {i n \theta}
\]

for all \(\theta\in\mathbf{R}\) .

\begin{enumerate}
\def\labelenumi{\alph{enumi})}
\setcounter{enumi}{4}
\item
  Let F be a subrepresentation of \(E_{n}\) . The orthogonal projector of \(E_{n}\) whose image is the \(\chi_{-n}\)-isotypic component of the restriction of \(\pi_{n}\) to \(\mathbf{SO}(\mathbf{R}^{2})\) is given by \(f \mapsto -(2i)^{n}f(i)f_{n}\) . (Use cor. 2 of V, p. 465 and the Cauchy integral formula, FVC, to appear.)
\item
  The representation \(\pi_n\) is irreducible.
\end{enumerate}

\(g)\) The matrix coefficient \(g \mapsto \langle f_n | \pi_n(g)f_n \rangle\) is nonzero and belongs to \(\mathcal{L}^2(\mathrm{G})\) .

\begin{enumerate}
\def\labelenumi{\alph{enumi})}
\setcounter{enumi}{7}
\tightlist
\item
  The representation \(\pi_{n}\) is square-integrable and there exists a constant c > 0 such that the formal degree of \(\pi_{n}\) is \(c(n - 1)\) .
\end{enumerate}

\begin{enumerate}
\def\labelenumi{\arabic{enumi})}
\setcounter{enumi}{16}
\item
  Let G be a connected real Lie group. There exists a faithful finite-dimensional unitary representation of G if and only if the group G is locally isomorphic to a compact group.
\item
  Let X be a locally compact space equipped with a bounded positive measure \(\mu\) . For every Banach space F, we denote by S(X, \(\mu\) ; F) the space of classes of functions \(\mu\) -measurable from X into F modulo the relation of equality \(\mu\) -almost everywhere, endowed with the topology of convergence in measure (INT, IV, p. 194, § 5, no. 11). We denote by A the commutative subalgebra of \(\mathcal{L}(\mathrm{L}_{\mathbf{C}}^{2}(\mathrm{X},\mu))\) formed by the multiplication endomorphisms by functions \(f \in \mathcal{L}_{\mathbf{C}}^{\infty}(\mathrm{X},\mu)\) .
\end{enumerate}

We denote by \(e\colon \mathbf{R}\to \mathbf{C}\) the homomorphism \(t\mapsto \exp (2i\pi t)\) .

Let E be a topological vector space.

\begin{enumerate}
\def\labelenumi{\alph{enumi})}
\tightlist
\item
  Let \(f \in \mathrm{S}(\mathrm{X}, \mu; \mathbf{R})\) such that \(|f(x)| \geqslant 1\) for \(\mu\) -almost all \(x \in \mathrm{X}\) . There exists a real number \(t \in ]0, 1[\) such that
\end{enumerate}

\[
\mu (\{x \in \mathrm{X} \mid t f (x) \in [ 1 / 4, 3 / 4 ] + \mathbf {Z} \}) \geqslant \frac {2}{5} \mu (\mathrm{X}).
\]

\begin{enumerate}
\def\labelenumi{\alph{enumi})}
\setcounter{enumi}{1}
\item
  Let u be a linear map from E into \(\mathrm{S}(\mathrm{X},\mu;\mathbf{R})\) . For every \(x\in E\) , we denote by \(\varrho_{u}(x)\) the element of A given by the endomorphism of multiplication by the function \(e\circ u(x)\) , which belongs to \(\mathcal{L}_{\mathbf{C}}^{\infty}(\mathrm{X},\mu)\) . The map \(\varrho_{u}\) is a continuous representation of G into \(\mathrm{L}_{\mathbf{C}}^{2}(\mathrm{X},\mu)\) if and only if the map u is continuous. (If u is not continuous, use the previous question to verify that \(\varrho\) is not continuous.)
\item
  Let \(\varrho\) be a unitary representation of E in \(\mathrm{L}^2 (\mathrm{X},\mu)\) whose image is contained in \(\mathcal{A}\) . There exists a unique continuous linear map \(u\colon \mathrm{E}\to \mathrm{S}(\mathrm{X},\mu ;\mathbf{R})\) such that \(\varrho = \varrho_{u}\) . (First consider the case where \(\mathrm{E} = \mathbf{R}\) , and then apply the Stone theorem.)
\end{enumerate}

\begin{enumerate}
\def\labelenumi{\arabic{enumi})}
\setcounter{enumi}{18}
\item
  Let G be a finite group and H a subgroup of G. For every unitary representation \(\pi\) of H in a finite-dimensional complex Hilbert space E, prove that the induced unitary representation \(\operatorname{Ind}_{\mathrm{H}}^{\mathrm{G}}(\pi)\) is isomorphic to the induced representation defined in A, VIII, p. 392, § 21, n° 3.
\item
  Let G be a locally compact group which is countable at infinity, and let H be a closed subgroup of G. Let \(\pi\) be a unitary representation of H in a Hilbert space E and let \(\varrho\) be the unitary representation of G induced by \(\pi\) . We denote by F the representation space of \(\varrho\) .
\end{enumerate}

We fix Haar measures \(\mu\) on G and \(\beta\) on H, and we denote by \(\kappa\colon\mathbf{G}\to\mathbf{R}_{+}^{*}\) a continuous function such that \(\kappa(xh)=\Delta_{\mathrm{H}}(h)\Delta_{\mathrm{G}}(h)^{-1}\kappa(x)\) for \((x,h)\in\mathbf{G}\times\mathrm{H}\) and such that \(\kappa(e)=1\) .

\begin{enumerate}
\def\labelenumi{\alph{enumi})}
\item
  If G/H admits a bounded G-invariant measure, and if there exists a nonzero vector \(x \in E\) which is H-invariant, then there exists a nonzero G-invariant vector \(y \in F\).
\item
  Conversely, suppose that \(y \in F\) is a nonzero and G-invariant vector of \(\varrho\) . Show that there exists a function \(f \in L_{\overline{\pi}}^{2}\) nonzero such that
\end{enumerate}

\[
f (x) = \left(\frac {\kappa (g ^ {- 1} x)}{\kappa (x)}\right) ^ {1 / 2} f (g ^ {- 1} x)
\]

for all \((g, x) \in \mathrm{G}^{2}\) . Deduce that \(\kappa = 1\) , and then the measure \(\mu/\beta\) on G/H is defined and bounded, and finally that \(f(e) \in \mathrm{E}\) is a nonzero H-invariant vector.
% semantic-fragments: legacy-input-seam
\relax{}
